%%%%%%%%%%%%%%%%%%%%%%%%%%%%%%%%%%%%%%%%%%%%%%%%%%%%%%%%%%%%%%%%%%%%%%%%%%%%%%%
%							Standard Paper Template
%							Simion Filip
%							Last updated 2026 Jan 24
%%%%%%%%%%%%%%%%%%%%%%%%%%%%%%%%%%%%%%%%%%%%%%%%%%%%%%%%%%%%%%%%%%%%%%%%%%%%%%%

\documentclass[11pt]{amsart}

\newif\ifdebug
% \debugtrue
\debugfalse

\newif \iffig
%\figtrue
\figfalse

\newif \iftable
%\tabletrue
\tablefalse


%%%%%%%%%%%%%%%%%%%%%%%%%%%%%%%%%%%%%%%%%%%%%%%%%%%%%%%%%%%%%%%%%%%%%
%		Packages

%==============================================================================
%							Standard AMS stuff
\usepackage{
	amsmath,
	amsfonts,
	amsthm,
	amssymb,
	bbm,
	mathrsfs,
	mathtools}
%==============================================================================
%	Link typesetting
% The hyphens option for {url} will allow breaks in long urls
\PassOptionsToPackage{hyphens}{url}\usepackage{hyperref}

% The code below alows further breaks
\makeatletter
\g@addto@macro{\UrlBreaks}{\UrlOrds}
\makeatother
%

% 	change \href colors
\usepackage{color}
\definecolor{darkred}{rgb}{0.4,0,0}
\definecolor{darkgreen}{rgb}{0,0.5,0}
\definecolor{darkblue}{rgb}{0,0,0.4}


%----------------------------------------------------
%			Autoref Fixes
%----------------------------------------------------
%	Some fixes are in the Theorem styles part
%
%	Section references
\renewcommand*{\sectionautorefname}{Section}

%	subsection \S gives paragraph symbol
%	the ~ binds them together
\def\subsek~{\S{}}
\renewcommand*\subsectionautorefname{\subsek}

\renewcommand*\subsubsectionautorefname{\subsek}

\renewcommand{\itemautorefname}{}


%	Makes \autoref do "Eqn. (3.1)"
%	instead of "Eqn. 3.1"
\def\equationautorefname~#1\null{%
  Eqn.~(#1)\null
}
%==============================================================================

%==============================================================================
%					Section styles and numbering
%==============================================================================

%
%	Documentation for \@startsection:
% The  \@startsection{<name>}{<level>}{<indent>}{<beforeskip>}{<afterskip>}{<style>}*[<altheading>]{<heading>}
% The part after the *, including the * is optional.
%
%   name: e.g., subsection
%   level: a number, denoting depth of section - e.g., chapter = 1, section = 2, etc.
%   indent: Indentation of heading from left margin
%   beforeskip: Absolute value = skip to leave above the heading.
% 	      If negative, then paragraph indent of text following heading is suppressed.
%   afterskip: If positive, then skip to leave below heading, else negative of skip to leave to right of run-in heading.
%   style: Commands to set style. Since June 1996 release the last command in this argument may be a command such as \MakeUppercase or  \fbox that takes an argument. The section heading will be supplied as the argument to this command. So setting #6 to, say, \bfseries\MakeUppercase would produce bold, uppercase headings.
%
% If * is missing, then increment the counter. If it is present, then there should be no [<altheading>] argument. The command uses the counter secnumdepth. It contains a pointer to the highest section level that is to be numbered.

%----------------------------------------------------
%
%	Hacking the (sub(sub))section commands
%
%
%	Boldness issue: I would like section and subsection not bold.
%	But I need the counter of subsubsections to be bold, and same for theorems.
%	This is \@secnumfont and it's one for all of these! (in amsart). Ridiculous.
%	So all counters have to be bold, so all titles have to be bold.
%
\makeatletter
\renewcommand{\@secnumfont}{\bfseries}
\makeatother
%		Redefining section
%	sections are centered, bold sans-serif
%  Format:
%{\@startsection{name}
% 				{1}   	% Level
% 				{\z@}	% ?
% 				{-3.5ex \@plus -1ex \@minus -.2ex} % Above spacing
%				{2.3ex \@plus.2ex} % After title spacing
%				{\centering \normalfont \Large \bfseries \sffamily}}% Title format
%	% the 5ex \@plus \@minus allow TeX to stretch more or less by those amounts

\makeatletter
\renewcommand{\section}
{\@startsection{section}
				   {1}
				   {\z@}%
                                   {-3.5ex \@plus -1ex \@minus -.2ex}%
                                   {2.3ex \@plus.2ex}%
                                   {\centering \normalfont \Large \bfseries \sffamily}}%
\makeatother

%		Redefining subsection
%	subsections are left-aligned, bold sans-serif
\makeatletter
\renewcommand{\subsection}
{\@startsection{subsection}
				     {2}
				     {\z@}%
                     {-3.25ex\@plus -1ex \@minus -.2ex}%
                     {1.5ex \@plus .2ex}%
                     {\normalfont \large \bfseries \sffamily}}% from \large
\makeatother

%		Redefining subsubsection
%	subsubsections are "2.3.1 Name" with Name in bold, no linebreak for subsequent text
\makeatletter
\renewcommand{\subsubsection}
{\@startsection{subsubsection}
				     {3}
				     {\z@}%
                     {1ex \@plus1ex \@minus.2ex}%
				     {-\fontdimen 2\font }
	                 {\normalfont\normalsize\bfseries}}
\makeatother

%----------------------------------------------------
%		Redefining paragraph
\makeatletter
\renewcommand{\paragraph}{%
\@startsection {paragraph}
				    {4}
				    {\z@}
				    {\z@}
				    {-\fontdimen 2\font }
				    {\normalfont\normalsize\bfseries}
}
\makeatother


%----------------------------------------------------
%		Enumerate and Itemize indices will be (i), (ii), etc.
%
\renewcommand{\labelenumi}{\theenumi}%
\renewcommand{\theenumi}{(\roman{enumi})}%
%----------------------------------------------------


%%----------------------------------------------------
%%		Enumerate Lists will be numbered
%%		automatically as (i), (ii) after
%%		just using \item
\renewcommand\labelenumi{(\roman{enumi})}
\renewcommand\theenumi\labelenumi
%==============================================================================






%==============================================================================
%								Typesetting
%==============================================================================
%
%----------------------------------------------------
%					Fonts
%----------------------------------------------------
%	{microtype} package will improve typesetting
%	[final] option is so that it also does it in
%	draft mode
\usepackage[stretch=10, final]{microtype}
%
%	For Latin Modern update to Computer Modern see
%	http://tex.stackexchange.com/questions/1390/latin-modern-vs-cm-super
\usepackage{lmodern}
%
%	[T1]{fontenc} adds more glyphs
%	Don't use with XeLaTeX; Just plain LaTeX
%	http://tex.stackexchange.com/questions/664/why-should-i-use-usepackaget1fontenc
%	Can lead to bad PDFs if not used with a modern font
%	http://tex.stackexchange.com/questions/1291/why-are-bitmap-fonts-used-automatically
\usepackage[T1]{fontenc}
%
%	Warning: fontenc + lmodern can sometimes create funny characters
%	especially if non-ASCII stuff appears.
%----------------------------------------------------
%						Date Fixing
%----------------------------------------------------
%	amsart places the date in a footnote
%	Patch below will place it below author name
%
\usepackage{etoolbox}
\makeatletter
\patchcmd{\@maketitle}
  {\ifx\@empty\@dedicatory}
  {\ifx\@empty\@date \else {\vskip3ex \centering\footnotesize\@date\par\vskip1ex}\fi
   \ifx\@empty\@dedicatory}
  {}{}
\patchcmd{\@adminfootnotes}
  {\ifx\@empty\@date\else \@footnotetext{\@setdate}\fi}
  {}{}{}
\makeatother
%----------------------------------------------------

% Define a command to apply the no-indent behavior to a specific environment
% requires etoolbox package
\makeatletter
\newcommand*\NoIndentAfterEnv[1]{%
  \AfterEndEnvironment{#1}{\par\@afterindentfalse\@afterheading}%
}
\makeatother

% Apply this to your desired environment, e.g., a custom 'myenv'
% \NoIndentAfterEnv{myenv}

% xspace is necessary for correct spacing after pre-defined names
\usepackage{xspace}

% recommended package for bold math
% https://tex.stackexchange.com/questions/595/how-can-i-get-bold-math-symbols
% works especially well for bold Greek or other symbols
% LATER: replace \mathbf{-} below by \bm{-}
\usepackage{bm}

% Keep using amsfonts
% for \mathbb
\usepackage{amsfonts}
\usepackage{amsmath}
\usepackage{amssymb}

\newcommand \cC {\mathcal{C}}
\newcommand \cV {\mathcal{V}}
\newcommand \cL {\mathcal{L}}
\newcommand \cG {\mathcal{G}}
\newcommand \bN {\mathbb{N}}
\newcommand \bR {\mathbb{R}}

% 	Inner product
\newcommand{\ip}[1]{\left\langle#1\right\rangle}
\newcommand{\Frechet}     {\xspace{{F}r\'echet}\xspace}
% epsilon
\newcommand{\ve}{\varepsilon}
%	approximate inequalities
\newcommand{\leqapprox}{\lesssim}
\newcommand{\geqapprox}{\gtrsim}
% Absolute value typeset nicely
\newcommand{\abs}[1]{\left|#1\right|}
% Norm of a vector
\newcommand{\norm}[1]{\left\|#1\right\|}
\newcommand{\SL}{\operatorname{SL}}
\newcommand{\id}{\mathbf{1}}


%%%%%%%%%%%%%%%%%%%%%%%%%%%%%%%%%%%%%%%%%%%%%%%%%%%%%%%%%%%%%%%%%%%%%%%%%%%%%%%
%                   Theorem styles and numbering
%%%%%%%%%%%%%%%%%%%%%%%%%%%%%%%%%%%%%%%%%%%%%%%%%%%%%%%%%%%%%%%%%%%%%%%%%%%%%%%
%

%--------------- New theorem style ------------------
%   spaces above and below are a bit larger
%   Keeps the . after theorem name
\newtheoremstyle{mytheoremstyle}% name
    {5pt}                       % Space above
    {5pt}                       % Space below
    {\itshape
    \hangindent \parindent}     % Body font - makes it indented by \parindent amount
    {}                          % Indent amount
    {\bfseries}                 % Theorem head font
    {.}                         % Punctuation after theorem head
    {.5em}                      % Space after theorem head
    {}              % Theorem head spec (can be left empty, meaning ‘normal’)

%--------------- New definition/etc style ------------------
%   spaces above and below are a bit larger (usually 3pt)
%   Keeps the . after definition name
%   Makes text normal.
\newtheoremstyle{mydefinitionstyle}% name
    {5pt}                       % Space above
    {5pt}                       % Space below
    {}                          % Body font
    {}                          % Indent amount
    {\bfseries}                 % Theorem head font
    {.}                         % Punctuation after theorem head
    {.5em}                      % Space after theorem head
    {}              % Theorem head spec (can be left empty, meaning ‘normal’)


% Switch to custom theorem style
\theoremstyle{mytheoremstyle}

%   package for fixing counters to work with autoref
\usepackage{aliascnt}

%%%%%%%%%%%%%%%%%%%%%%%%%%%%%%%%%%%%%%%%%%%%%%%%%%%%%%%%%%%%%%%%%%%%%%%%%%%%%%%
%                   Theorems for Introduction
%%%%%%%%%%%%%%%%%%%%%%%%%%%%%%%%%%%%%%%%%%%%%%%%%%%%%%%%%%%%%%%%%%%%%%%%%%%%%%%

%   Main Theorems are in the introduction.
%   Will be labeled Theorem A, Theorem B, etc.
\newtheorem{theoremmain}{Theorem}
\renewcommand{\thetheoremmain}{\Alph{theoremmain}}

\providecommand*{\theoremmainautorefname}{The\-o\-rem}


% Theorem in the introduction are numbered independently (1,2,3 etc)
% add [section] to have 1.1, 1.2 etc.
\newtheorem{theoremintro}{Theorem}
\providecommand*{\theoremintroautorefname}{The\-o\-rem}



% INTRO CONJECTURES are numbered as intro theorems
\newaliascnt{conjectureintro}{theoremintro}
\newtheorem{conjectureintro}[conjectureintro]{Conjecture}
\aliascntresetthe{conjectureintro} 
\providecommand*{\conjectureintroautorefname}{Con\-jec\-ture}

% INTRO COROLLARIES are numbered as intro theorems
\newaliascnt{corollaryintro}{theoremintro}
\newtheorem{corollaryintro}[corollaryintro]{Corollary}
\aliascntresetthe{corollaryintro} 
\providecommand*{\corollaryintroautorefname}{Co\-rollary}


% making INTRO REMARKS and QUESTIONS plain text style
\theoremstyle{mydefinitionstyle}

% INTRO QUESTIONS are numbered as intro theorems
\newaliascnt{questionintro}{theoremintro}
\newtheorem{questionintro}[questionintro]{Question}
\aliascntresetthe{questionintro} 
\providecommand*{\questionintroautorefname}{Ques\-tion}

% INTRO REMARKS are numbered as intro theorems
\newaliascnt{remarkintro}{theoremintro}
\newtheorem{remarkintro}[remarkintro]{Remark}
\aliascntresetthe{remarkintro} 
\providecommand*{\corollaryintroautorefname}{Re\-mark}


\theoremstyle{mytheoremstyle}


%%%%%%%%%%%%%%%%%%%%%%%%%%%%%%%%%%%%%%%%%%%%%%%%%%%%%%%%%%%%%%%%%%%%%%%%%%%%%%%
%                  Theorems and environments for main body
%%%%%%%%%%%%%%%%%%%%%%%%%%%%%%%%%%%%%%%%%%%%%%%%%%%%%%%%%%%%%%%%%%%%%%%%%%%%%%%
% change numbering to "2.3.1 Theorem" rather than "Theorem 2.3.1"
\swapnumbers

% Theorem numbers are within SECTIONS
\newtheorem{theorem}{Theorem}[section]

%   Hyphenate Theorem references
\renewcommand{\theoremautorefname}{The\-o\-rem}

% Equations are numbered within SECTIONS
\numberwithin{equation}{section}

% Makes theorems numbered with the same counter as equations
\makeatletter
    \let\c@theorem \c@equation
\makeatother

% Makes subsections numbered with the same counter as equations
\makeatletter
    \let\c@subsection\c@equation
\makeatother

%       Redefining subsection
%   subsections are left-aligned, bold sans-serif
\makeatletter
\renewcommand{\subsection}
{\@startsection{subsection}
                     {2}
                     {\z@}%
                     {1ex \@plus1ex \@minus.2ex}%
                     {-\fontdimen 2\font }
                     {\normalfont\normalsize\bfseries}}
\makeatother


%       NEXT INSTRUCTIONS ARE NOW IN package_settings.sty
% Makes figures numbered with the same counter as equations
% \makeatletter
%   \let\c@figure\c@equation      % Note that need to include (above) \numberwithin{figure}{subsubsection}
% \makeatother


%%%%%%%%%%%%%%%%%%%%%%%%%%%%%%%%%%%%%%%%%%%%%%%%%%%%%%%%%%%%%%%%%%%%%%%%%%%%%%%
%       Creates Lemma, Proposition, etc. environments
%
%   Fixes \autoref numberings, creates alias counters
%   Now \autoref will produce Lemma 2.3.1 etc.
%
%   based on aliascnt package


\newaliascnt{conjecture}{theorem}  
\newtheorem{conjecture}[conjecture]{Conjecture}
\aliascntresetthe{conjecture} 
\providecommand*{\conjectureautorefname}{Con\-jec\-ture}

\newaliascnt{lemma}{theorem}  
\newtheorem{lemma}[lemma]{Lemma}  
\aliascntresetthe{lemma} 
\providecommand*{\lemmaautorefname}{Lem\-ma}

\newaliascnt{proposition}{theorem}
\newtheorem{proposition}[proposition]{Proposition}
\aliascntresetthe{proposition} 
\providecommand*{\propositionautorefname}{Pro\-po\-si\-tion}

\newaliascnt{corollary}{theorem}  
\newtheorem{corollary}[corollary]{Corollary}  
\aliascntresetthe{corollary} 
\providecommand*{\corollaryautorefname}{Co\-ro\-lla\-ry}



%%%%%%%%%%%%%%%%%%%%%%%%%%%%%%%%%%%%%%%%%%%%%%%%%%%%%%%%%%%%%%%%%%%%%%%%%%%%%%%
%       Changes to custom definition style (normal body text)
\theoremstyle{mydefinitionstyle}



\newaliascnt{question}{theorem}  
\newtheorem{question}[question]{Question}  
\aliascntresetthe{question} 
\providecommand*{\questionautorefname}{Ques\-tion}


\newaliascnt{problem}{theorem}  
\newtheorem{problem}[problem]{Problem}  
\aliascntresetthe{problem} 
\providecommand*{\problemautorefname}{Prob\-lem}


\newaliascnt{exercise}{theorem}  
\newtheorem{exercise}[exercise]{Exercise}  
\aliascntresetthe{exercise} 
\providecommand*{\exerciseautorefname}{Ex\-er\-cise}

\newaliascnt{definition}{theorem}  
\newtheorem{definition}[definition]{Definition}  
\aliascntresetthe{definition} 
\providecommand*{\definitionautorefname}{De\-fi\-ni\-ti\-on}

\newaliascnt{remark}{theorem}  
\newtheorem{remark}[remark]{Remark}  
\aliascntresetthe{remark} 
\providecommand*{\remarkautorefname}{Re\-mark}

\newaliascnt{example}{theorem}  
\newtheorem{example}[example]{Example}  
\aliascntresetthe{example} 
\providecommand*{\exampleautorefname}{Ex\-am\-ple}



\NoIndentAfterEnv{theorem}
\NoIndentAfterEnv{proposition}
\NoIndentAfterEnv{corollary}
\NoIndentAfterEnv{lemma}
\NoIndentAfterEnv{problem}
\NoIndentAfterEnv{question}


%%%%%%%%%%%%%%%%%%%%%%%%%%%%%%%%%%%%%%%%%%%%%%%%%%%%%%%%%%%%%%%%%%%%%
%				Title & Date
%%%%%%%%%%%%%%%%%%%%%%%%%%%%%%%%%%%%%%%%%%%%%%%%%%%%%%%%%%%%%%%%%%%%%
%	also: Add it in the hyperref setup
%
\title[Zero entropy for globally hypoelliptic vector fields]{Zero entropy for globally hypoelliptic vector fields}
%	change to \title[short title]{Actual title}
%	if what appears on headers is too long

%	Hyperref options
\hypersetup{
	pdftitle={Zero entropy for globally hypoelliptic vector fields},	% title
	pdfauthor={Simion Filip},	% author
	pdfsubject={Zero entropy for globally hypoelliptic vector fields},		% subject of the document
	pdfkeywords={MSC},	% list of keywords
	pdfnewwindow=true,		% links in new window
	colorlinks=true,		% false: boxed links; true: colored links
	linkcolor=darkblue,		% color of internal links (change box color with linkbordercolor)
	citecolor=darkred,		% color of links to bibliography
	filecolor=darkblue,		% color of file links
	urlcolor=darkblue,		% color of external links
	pdfborder={0 0 0},
	breaklinks=true
}


%---------------------------------------
%						Date of revision
\thanks{Revised \textsc{\today} }

%						Date First posted
\date{August 2026}

%%%%%%%%%%%%%%%%%%%%%%%%%%%%%%%%%%%%%%%%%%%%%%%%%%%%%%%%%%%%%%%%%%%%%
%						Author info
%%%%%%%%%%%%%%%%%%%%%%%%%%%%%%%%%%%%%%%%%%%%%%%%%%%%%%%%%%%%%%%%%%%%%

\author{
	Simion Filip
}
%


%%%%%%%%%%%%%%%%%%%%%%%%%%%%%%%%%%%%%%%%%%%%%%%%%%%%%%%%%%%%%%%%%%%%%%%%%%%%%%%
%%%%%%%%%%%%%%%%%%%%%%%%%%%%%%%%%%%%%%%%%%%%%%%%%%%%%%%%%%%%%%%%%%%%%%%%%%%%%%%
%							Start of Main Document
%%%%%%%%%%%%%%%%%%%%%%%%%%%%%%%%%%%%%%%%%%%%%%%%%%%%%%%%%%%%%%%%%%%%%%%%%%%%%%%
%%%%%%%%%%%%%%%%%%%%%%%%%%%%%%%%%%%%%%%%%%%%%%%%%%%%%%%%%%%%%%%%%%%%%%%%%%%%%%%

\begin{document}
\begin{center}
{\large Simion Filip}
\end{center}
\section{Introduction}
	\label{sec:introduction}


Let $M$ be a smooth compact connected manifold, orientable and without boundary, of dimension $m$.
Denote by $\cC^\infty(M)$ the Fr\'echet space of smooth functions on $M$ and by $\cV^\infty(M)$ the Fr\'echet space of smooth top degree forms on $M$.
Denote by $\cC^{-\infty}(M)$ the topological dual of $\cV^\infty(M)$, whose elements will be called generalized functions on $M$, and by $\cV^{-\infty}(M)$ the topological dual of $\cC^\infty(M)$, whose elements will be called generalized measures on $M$.
Note that measures belong to $\cV^{-\infty}(M)$.

Let now $X$ be a smooth vector field on $M$ and denote by $\cL_X$ the Lie derivative along $X$.

\begin{definition}[Globally hypoelliptic/Cohomology-free vector field]
	\label{def:Hypoelliptic-Cohomology-free-vector-field}
	The vector field $X$ is called \emph{globally hypoelliptic} if, whenever $u$ is a generalized function on $M$ such that $\cL_Xu$ is smooth, we have that $u$ is smooth.

	The vector field $X$ is called \emph{cohomology-free} if the map
	\[
	\cL_X:\cC^\infty(M)\to\cC^\infty(M)
	\]
	has closed range and its cokernel is one-dimensional.
\end{definition}
A vector field is globally hypoelliptic if and only if it is cohomology-free, by \cite[Thm.~3.10]{Forni2008_On-the-Greenfield-Wallach-and-Katok-conjectures-in-dimension-three} (extending the proof on the torus by Chen--Chi \cite[Thm.~2.2]{ChenChi2000_Hypoelliptic-vector-fields-and-almost-periodic-motions-on-the-torus}, see also \cite[\S4]{Rodriguez-HertzRodriguez-Hertz2006_Cohomology-free-systems-and-the-first-Betti-number}).
We will refer to a vector field satisfying either of the equivalent assumptions in \autoref{def:Hypoelliptic-Cohomology-free-vector-field} as a $KGW$-vector field, for \cite[Conj.~3.6]{Katok2001_Cocycles-cohomology-and-combinatorial-constructions-in-ergodic-theory} (who called them cohomology-free) and Greenfield--Wallach \cite[Problem~2]{GreenfieldWallach1973_Globally-hypoelliptic-vector-fields} who called them globally hypoelliptic.
In each of these works, it was conjectured that the only KGW vector fields are linear Diophantine ones on tori.

KGW vector fields satisfy a number of known properties, including the following:
\begin{enumerate}
	\item The vector space of invariant generalized measures is one-dimensional, spanned by a smooth positive fully supported volume form on $M$ normalized to $\ip{\mu,1}=1$.
	\item On functions with zero average for $\mu$, denoted $\cC^\infty_0(M)$, the operator $\cL_X$ is an isomorphism of \Frechet spaces.
\end{enumerate}
Our main result is:
\begin{theorem}[KGW has zero entropy]
	\label{thm:KGW-has-zero-entropy}
	The flow generated by a $KGW$-vector field has zero topological entropy.
\end{theorem}

Our proof is based on several ingredients, of possibly independent interest, which we now describe.

The first is a general property of dynamical systems with positive entropy on manifolds; as will be clear from the proof, the result can be generalized to other spaces with polynomial growth of the number of balls of a given radius needed to cover the space.


\begin{lemma}[Exponential recurrence]
	\label{lemma:Exponential-recurrence-flows}
	Let $g_t\colon M\to M$ be a smooth flow on a compact manifold, with positive topological entropy and no fixed points.

	There exist $c,\ve>0$, a sequence of points $x_i,y_i\in M$, and times $T_i\to +\infty$, satisfying:
	\begin{enumerate}
		\item $d(x_i,y_i)\leq e^{-\ve T_i}$
		\item $d\left(g_{T_i}x_i,g_{T_i}y_i\right)\leq e^{-\ve T_i}$
		\item 
		\label{item:exp_rec_separation_flows}
		For any continuous strictly increasing $\rho\colon [0,T_i]\to [0,T_i]$ with $\rho(0)=0,\rho(T_i)=T_i$, there exists $t_{i,\rho}\in [0,T_i]$ such that $d\left(g_{\rho(t_{i,\rho})}x_i,g_{t_{i,\rho}}y_i\right)\geq c$.
	\end{enumerate}
\end{lemma}

The second result we will need applies to ergodic integrals of the flow $g_t$ generated by a KGW vector field $X$.
Denote by $\delta_{[x,g_tx]}$ the measure on $M$ supported on the corresponding orbit, defined explicitly on functions by
\[
	\ip{\delta_{[x,g_tx]},f}:=\int_{0}^t f(g_s x)ds.
\]
We then have:
\begin{lemma}[Endpoint distributions]
	\label{lemma:Endpoint-distributions}
	Suppose $X$ is a KGW vector field.
	Then there exists a family of generalized measures $\xi_x\in \cV^{-\infty}(M)$, defined for every $x\in M$, with the following properties:
	\begin{enumerate}
		\item We have the identity of generalized measures:
		\begin{align}
			\label{eqn:distribution-orbit-identity}
			\delta_{[x,g_tx]} = t \mu + (\xi_{g_tx}-\xi_x)
		\end{align}
		\item We have the Lipschitz bound:
		\begin{align}
			\label{eqn:Lipschitz-bound-xi}
			\norm{\xi_x-\xi_y}_{\cV^{-k_0}} \leqapprox d(x,y)
			\text{ and }
			\quad \norm{\xi_x}_{\cV^{-k_0}} \leqapprox 1
			\quad \forall x,y\in M
		\end{align}
		for some $k_0\in \bN$, where $\cV^{-k_0}$ denotes the norm on generalized measures obtained by testing against $\cC^{k_0}$-functions.
	\end{enumerate}
\end{lemma}

The last result has an elementary but useful consequence:

\begin{corollary}[Diophantine properties]
	\label{cor:diophantine_properties}
	Suppose $X$ is a KGW vector field, and $\dim M\geq 2$.
	Then there exist $C,c,A,a>0$ such that, for any $x\in M$ and $T\geq 1$ we have:
	\begin{align}
		\label{eqn:Diophantine_repulsion}
		d(x,g_Tx) \geq \frac{c}{T^{A}}
	\end{align}
	and for any $y\in M$ there exists $t\in [0,T]$ such that
	\begin{align}
		\label{eqn:Diophantine_equidistribution}
		d(y,g_t x) \leq \frac{C}{T^{a}}.
	\end{align}	
\end{corollary}
Note that the Diophantine properties in this corollary are the conjectured ones in the case of a linear flow on the torus.
At the same time, the horocycle flow on a compact quotient of $\SL_2(\bR)$ satisfies the conclusions of this corollary.

The proofs of the above results are in the next section.

%%=============================================================================
%%          start of subsec: Discrete-time case

\subsection{Discrete-time case}
  \label{ssec:discrete_time_case}

We note that the main result \autoref{thm:KGW-has-zero-entropy} generalizes straightforwardly to the discrete-time case.
Namely, a diffeomorphism $f\colon M\to M$ is \emph{cohomology-free} if the operator $f^*-\operatorname{id}$ acting on $\cC^\infty(M)$ has image a closed subspace of codimension $1$.
Then such a diffeomorphism satisfies the analogues of all the above stated results.
The proofs are essentially the same, with the operator $\cL_X$ replaced by $f^*-\operatorname{id}$.
\autoref{lemma:Exponential-recurrence-flows} becomes in fact easier, as one can appeal to the Brin--Katok local entropy formula \cite{BrinKatok1983_On-local-entropy}, one does not need to consider reparametrized orbits, and the matching \autoref{lem:orbit_matching} below is straightforward for discrete time dynamics.

%%          end of subsec: Discrete-time case
%%=============================================================================



%%%%%%%%%%%%%%%%%%%%%%%%%%%%%%%%%%%%%%%%%%%%%%%%%%%%%%%%%%%%%%%%%%%%%%%%%%%%%%%
% 						End of Section: Introduction
%%%%%%%%%%%%%%%%%%%%%%%%%%%%%%%%%%%%%%%%%%%%%%%%%%%%%%%%%%%%%%%%%%%%%%%%%%%%%%%



%%%%%%%%%%%%%%%%%%%%%%%%%%%%%%%%%%%%%%%%%%%%%%%%%%%%%%%%%%%%%%%%%%%%%%%%%%%%%%%
%%%%%%%%%%%%%%%%%%%% Section Proofs %%%%%%%%%%%%%%%%%%%%
%%%%%%%%%%%%%%%%%%%%%%%%%%%%%%%%%%%%%%%%%%%%%%%%%%%%%%%%%%%%%%%%%%%%%%%%%%%%%%%

					\section{Proofs}
					\label{sec:Proofs}

We will use the notation $A\leqapprox B$ when there exists a constant $c>0$ such that $A\leq cB$, where $c$ is allowed to depend on the manifold and vector field and perhaps other a priori fixed choices (e.g. charts, Riemannian metric, etc.).
We will also write $A\asymp B$ for $A\leqapprox B$ and $B\leqapprox A$.


%%=============================================================================
%%          start of subsec: Proof of exponential recurrence flows

\subsection{Proof of \autoref{lemma:Exponential-recurrence-flows}}
  \label{ssec:proof_of_exponential_recurrence_flows}
We divide the proof into two natural parts.

\noindent\textbf{Step 1: Entropy considerations.}
Consider the set of reparametrizations:
\[
  Rp[0,T] =\{\rho\colon [0,T]\to [0,T]\text{ orientation-preserving homeomorphism}\}.
\]
Recall that ordinary Bowen balls are defined by
\[
  B^T(x,r):=\{y\in M\colon d(g_t x, g_t y)<r, \quad \forall t\in [0,T]\}
\]
and consider the Bowen balls with ``reparametrization'':
\[
  B^T_{rp}(x,r):= \{y\in M\colon \exists \rho\in Rp[0,T] \text{ s.t. }d(g_{\rho(t)} x, g_{t}y)< r
  ,\forall t\in [0,T]\}
\]
We note for future reference that
\begin{align}
  \label{eqn:reparam_ball_symmetry}
  y\in B^T_{rp}(x,r) \iff x\in B^T_{rp}(y,r)  
\end{align}
by considering the inverse parametrization, and similarly:
\begin{align}
  \label{eqn:reparam_ball_triangle}
  y\in B^T_{rp}(x,r_1) \text{ and } z \in B^T_{rp}(y,r_2)
  \implies 
  z \in B^T_{rp}(x,r_1+r_2).
\end{align}

Denote by $N^T_{rp}(r)$ the minimal number of $x_i\in M$ such that $B^T_{rp}(x_i,r)$ cover $M$, and $N^T(r)$ for the minimal number using ordinary Bowen balls.
Set
\[
  h_{rp}(g_\bullet):=\lim_{r\to 0}\limsup_{T\to +\infty} \frac{1}{T}\log N^T_{rp}(r)
\]
and analogously $h_{top}(g_\bullet)$, which is the topological entropy of the flow, which is also the topological entropy of the diffeomorphism $g_1$.

A result of Thomas \cite[Thm.~10]{Thomas1990_Topological-entropy-of-fixed-point-free-flows} implies that $h_{rp}(g_\bullet)=h_{top}(g_\bullet)>0$.
Indeed, it is clear that the reparametrized Bowen balls are at least as large as ordinary Bowen balls, and the balls used by Thomas in loc. cit. are larger than the $B_{rp}^T(x,r)$ defined above (Thomas uses a larger class of reparametrizations).

Let $E^T(r)\subset M$ denote a set of minimal cardinality such that $B^T_{rp}(x_i,r)$ cover $M$.
It follows that there exists $r>0$ and a sequence $T_i\to +\infty$ such that we have
\[
  \abs{E^{T_i}(r)}=N^{T_i}_{rp}(r)\geq e^{T_i\beta} \quad \text{ for }\beta:=\tfrac 12h_{top}(g_\bullet)>0.
\]
Note that $N^T_{rp}(r')\geq N^T_{rp}(r)$ if $r'<r$.
We work with $T$ belonging to the sequence $T_i$ just obtained.

Consider now a subset $\hat{E}^T(r/4)$ of $E^T(r/4)$, of maximal cardinality with the property that the balls $B^T_{rp}(\hat{e},r/4)$ over all $\hat{e}\in \hat{E}^T(r/4)$ are disjoint.
Then we claim that the union of all the balls $B^T_{rp}(\hat{e},r)$ over all $\hat{e}\in \hat{E}^T(r/4)$ cover $M$ and hence
\begin{align}
  \label{eqn:exponential_growth_reparam_balls}
  \abs{\hat{E}^T(r/4)} \geq \abs{E^T(r)}\geq e^{T\beta}.
\end{align}
Indeed, for any $z\in M$ there exists $y\in E^T(r/4)$ such that $z\in B^T_{rp}(y,r/4)$.
If $y\in \hat{E}^T(r/4)$ we are done, if not $y$ cannot be added in so there exists $x\in \hat{E}^T(r/4)$ such that $B^T_{rp}(x,r/4)\cap B^T_{rp}(y,r/4)\neq \emptyset $.
It follows from \autoref{eqn:reparam_ball_symmetry} and \autoref{eqn:reparam_ball_triangle} that $z\in B^T_{rp}(x,r)$.


\noindent\textbf{Step 2: Polynomial pigeonhole.}
Since $M$ is a compact manifold of dimension $m$, there exists a constant $C\geq 1$ such that for any $\delta>0$, there exists a set $P_\delta\subset M$ of cardinality at most $C\delta^{-m}$ such that every point in $M$ is at (ordinary) distance at most $\delta$ to some element of $P_\delta$.
Set $\delta=\delta_T:=e^{-T\tfrac{\beta}{4m}}$, so $T$ remains the only free parameter.
Set also $E^T:=\hat{E}^T(r/4)$ satisfying the exponential lower bound in \autoref{eqn:exponential_growth_reparam_balls}.

Choose now maps $\alpha,\omega\colon E^T\to P_\delta$ such that
\begin{align*}
  d(s,\alpha(s)) & \leq \delta \quad \forall s\in E^T\\
  d(g_T s,\omega(s)) & \leq \delta \quad \forall s\in E^T
\end{align*}
which is possible by the properties of $P_\delta$.
Now the cardinality of $P_\delta\times P_\delta$ is at most $C^2 e^{T\beta/2}$, while the cardinality of $E^T$ is at least $e^{T\beta}$ by \autoref{eqn:exponential_growth_reparam_balls}.
Assuming $T\geq T_0$ for $T_0$ sufficiently large, the map $\alpha\times \omega\colon E^T \to P_\delta\times P_\delta$ cannot be injective.
Thus, there exist distinct $s_1,s_2\in E^T$ with the following properties:
\begin{enumerate}
  \item $d(s_1,s_2)\leq 2e^{-T \tfrac{\beta}{4m}}$ since $\alpha(s_1)=\alpha(s_2)$.
  \item $d(g_Ts_1,g_Ts_2)\leq 2e^{-T \tfrac{\beta}{4m}}$ since $\omega(s_1)=\omega(s_2)$.
  \item The balls $B^T_{rp}(s_1,r/4)$ and $B^T_{rp}(s_2,r/4)$ are disjoint, in particular $s_2\notin B^T_{rp}(s_1,r/4)$.
\end{enumerate}
Passing now to the sequence $T=T_i\to +\infty$ completes the proof of \autoref{lemma:Exponential-recurrence-flows}. \hfill \qed


%%          end of subsec: Proof of exponential recurrence flows
%%=============================================================================


%%=============================================================================
%%					start of subsec: Proof of endpoint distributions

\subsection{Proof of \autoref{lemma:Endpoint-distributions}}
	\label{ssec:proof_of_endpoint_distributions}
We know that $\cL_X\colon \cC^\infty_0(M)\to \cC^\infty_0(M)$ is an isomorphism of \Frechet spaces.
Denote by $\cG$ its inverse, extended to all smooth functions by the condition $\cG(1):=1$, which is a continuous isomorphism of $\cC^\infty(M)$.
Therefore, its dual is a continuous isomorphism
\[
	\cG^t \colon \cV^{-\infty}(M) \to \cV^{-\infty}(M).
\]
Note that we can decompose any generalized measure $\xi$ as $\xi = \xi_0 + \ip{\xi,1}\mu$ where $\ip{\xi_0,1}=0$.
Taking into account the normalization $\cG(1)=1$, we have that $\cG^t(\mu)=\mu$ and furthermore
\[
	\cG^t(\xi) = \cG^t(\xi_0) + \ip{\xi,1}\mu
\]
with the property that $\ip{\cG^t(\xi_0),1}=0$.

The continuity of $\cG$ implies that there exists $k_0\in \bN$ such that
\[
	\cG\colon \cC^{k_0}(M) \to \cC^1(M) \text{ is bounded}.
\]
Therefore
\begin{align}
	\label{eqn:Gt_bounded_C1}
	\cG^t\colon \cV^{-1}(M) \to \cV^{-k_0}(M) \text{ is bounded}.
\end{align}
For a point $x\in M$ denote by $\delta_x\in \cV^{-\infty}(M)$ the Dirac-delta measure supported at $x$.
Then $\delta_x\in \cV^0(M)$, i.e. it defines a continuous linear functional on continuous functions, and in the weaker $\cV^{-1}$-norm we have:
\[
	\norm{\delta_x-\delta_y}_{\cV^{-1}} \leq d(x,y)
\]
which follows from
\[
	|\ip{\delta_x-\delta_y,\phi}|\leq d(x,y)\norm{\phi}_{\cC^1}
\]
for any once-differentiable $\phi$.
Set
\begin{align}
	\label{eqn:xi_x_def_and_property}
	\xi_x:=\cG^t(\delta_x) \text{ so that }
	\cL_X(\xi_x) = \mu - \delta_x
\end{align}
More explicitly, for a smooth test function $\phi$ we have $\ip{\xi_x,\phi}=u(x)$ where $u$ is the unique solution to $\cL_X u = \phi - \int\! \phi d\mu$ with $\int\! ud\mu = 0$.
To justify \autoref{eqn:xi_x_def_and_property} note that $\cG \cL_X \phi = \phi - \ip{\mu,\phi}$ for any smooth test function $\phi$ and therefore
\begin{align*}
	\ip{\cL_X \cG^t \delta_x, \phi} & = 
	- \ip{\cG^t\delta_x, \cL_X \phi}
	= -\ip{\delta_x,\cG \cL_X\phi}\\
	& = 
	-\ip{\delta_x, \phi - \ip{\mu,\phi}}
	= \ip{\mu - \delta_x, \phi}.
\end{align*}
Note the estimate
\[
	\norm{\xi_x-\xi_y}_{\cV^{-k_0}}
	=
	\norm{\cG^t(\delta_x-\delta_y)}_{\cV^{-k_0}}
	\leqapprox
	\norm{\delta_x-\delta_y}_{\cV^{-1}}
	\leq
	d(x,y)
\]
where the middle inequality follows from \autoref{eqn:Gt_bounded_C1}, which implies the second bound in \autoref{eqn:Lipschitz-bound-xi} as well.

The claimed identity \autoref{eqn:distribution-orbit-identity}
\[
	\delta_{[x,g_tx]} = t \mu + (\xi_{g_tx}-\xi_x)
\]
follows by observing that both sides agree after applying $\cL_X$, hence the difference is $\cL_X$-invariant so a multiple of $\mu$, and
both sides yield $t$ when paired against the constant function $1$, so must agree.
Indeed $\cL_X(\delta_{[x,g_tx]}) =  \delta_x - \delta_{g_t x}$, since
\begin{align*}
	\ip{\cL_{X}\delta_{[x,g_tx]},\phi} & = 
	- \ip{\delta_{[x,g_t x]}, \cL_X \phi}
	= - \int_{0}^t (\cL_X\phi)(g_s x)ds\\
	& = 
	- [\phi(g_tx) - \phi(x)]
	=
	\ip{\delta_x - \delta_{g_t x},\phi}
\end{align*}
and $\cL_X(\xi_{g_t x} -\xi_x) = \delta_x - \delta_{g_tx}$ by \autoref{eqn:xi_x_def_and_property}.
This completes the proof of \autoref{lemma:Endpoint-distributions}.
\hfill \qed


%%					end of subsec: Proof of endpoint distributions
%%=============================================================================

%%=============================================================================
%%					start of subsec: Preliminary constructions in charts

\subsection{Preliminary constructions in charts}
	\label{ssec:preliminary_constructions_in_charts}
For the proofs of the remaining statements, we will exploit the estimates on the generalized measures $\xi_x$ obtained in \autoref{lemma:Endpoint-distributions} by testing them against bump functions concentrated at polynomially small scales and centered at appropriately selected points in the manifold.
To do so, we need some preliminaries.

Fix finitely many coordinate charts covering $M$, with images convex sets in $\bR^m$.
For each point of $M$ choose one of the charts such that the distance (in $\bR^m$) from the point to the boundary of the chart is uniformly bounded below by some $r_{ch}>0$.

We also assume that in these charts, the vector field is ``rectified'' to be equal to $\frac{\partial}{\partial x_1}$ where $x_1$ is the first coordinate.
This rectification is possible since the vector field has no fixed points, since the only invariant measure is volume.
Finally, we will measure distances in these charts using Euclidean distances, which yield estimates for the distance in the manifold up to uniform factors.

Fix a smooth nonnegative bump function $\rho$ on $\bR^m$, with support in the ball of radius $1$, and equal to $1$ on the ball of radius $\tfrac 12$.
Denote by $\rho_{r}(x):=\rho(x/r)$ and, for $p\in M$ denote by $\rho_{r,p}\in \cC^\infty(M)$ the placement of $\rho_r$ in $M$, using the chart associated to $p$, and centered at $p$.
We will only take $r>0$ satisfying $r\in(0,r_{ch})$.
We record the following properties of the functions:
\begin{align}
	\label{eqn:rho_rp_Ck_norm}
	\norm{\rho_{r,p}}_{\cC^{k_0}} & \leqapprox r^{-k_0}\\
	\ip{\mu,\rho_{r,p}}  & \asymp r^{m}
	\label{eqn:rho_rp_integral}
\end{align}


%%					end of subsec: Preliminary constructions in charts
%%=============================================================================

%%=============================================================================
%%					start of subsec: Proof of Diophantine corollary

\subsection{Proof of \autoref{cor:diophantine_properties}}
	\label{ssec:proof_of_diophantine_corollary}
To establish \autoref{eqn:Diophantine_equidistribution}, we will consider the bump function $\rho_{r,y}$ for some $y\in M$, and pair it against the distributional identity from \autoref{eqn:distribution-orbit-identity}:
\begin{align}
  \label{eqn:xi_x_bump_paired}
  \ip{\delta_{[x,g_Tx]}, \rho_{r,y} } = \ip{T \mu + (\xi_{g_Tx} - \xi_x), \rho_{r,y}}.
\end{align}
For any $y\in M$ we have
\[
  \ip{T\mu,\rho_{r,y}}\asymp Tr^{m} \text{ and }
  |\ip{\xi_{g_Tx}-\xi_{x},\rho_{r,y}}|
  \leqapprox r^{-k_0}.
\]
where we used that $\norm{\xi_{p}}_{\cV^{-k_0}}\leqapprox 1$ for all $p\in M$.
Select now $r:=T^{-\beta}$, so the first term is of the order $T^{1-m\beta}$ while the second is bounded above by $T^{\beta k_0}$.
Choosing $\beta>0$ sufficiently small (e.g. $\beta:=\tfrac{1}{2(m+k_0)}$) guarantees that the right-hand side of \autoref{eqn:xi_x_bump_paired} is nonzero.
Therefore, the orbit segment $[x,g_Tx]$ must pass through the $CT^{-\beta}$-neighborhood of $y$, for some $C>0$ sufficiently large but fixed.

To establish \autoref{eqn:Diophantine_repulsion}, we note that there exists a constant $b>0$ such that for any orbit segment of length $T$, there exists some $y\in M$ which is at distance at least $\tfrac{b}{T^{1/(m-1)}}$ from any point in the orbit.
This is valid for any smooth flow: the $R$-neighborhood of the orbit segment has volume at most $B(R^m+R^{m-1}T)$ for some $B>0$, by considerations in local charts: cover an orbit of length $T$ by $O(1 + T/R)$ balls of radius $R$.
By choosing $R=\tfrac{b}{T^{1/(m-1)}}$ for sufficiently small $b>0$, some point in $M$ will not be covered.

Returning to \autoref{eqn:distribution-orbit-identity}, choose such a $y\in M$ and $r:={b}T^{-\gamma}$ with $\gamma:=\tfrac{1}{m-1}$.
Then $\ip{\delta_{[x,g_Tx]},\rho_{r,y}}=0$.
At the same time
\[
	|\ip{\xi_{g_T x} - \xi_x,\rho_{r,y}}|\leqapprox d(x,g_Tx) T^{k_0\gamma}
\]
while $\ip{T\mu,\rho_{r,y}}\asymp T^{1-\gamma m}$.
It then follows that
\[
	d(x,g_Tx)\geqapprox T^{1-\gamma(m+k_0)} = T^{-\tfrac{k_0+1}{m-1}}
\]
and this concludes the proof. \hfill \qed

%%					end of subsec: Proof of Diophantine corollary
%%=============================================================================


%%=============================================================================
%%					start of subsec: Proof of zero entropy

\subsection{Proof of \autoref{thm:KGW-has-zero-entropy}}
	\label{ssec:proof_of_zero_entropy}
Suppose, by contradiction, that $g_t\colon M\to M$ has positive topological entropy.
Note that the case of $\dim M\leq 1$ is immediate.
We will derive a contradiction to \autoref{lemma:Exponential-recurrence-flows}\autoref{item:exp_rec_separation_flows}, so assume its conclusion holds and let $x_i,y_i, T_i$ be the corresponding points and times.


Set $r:=bT_i^{-\gamma}$ and select $b>0$ sufficiently small, $\gamma>0$ sufficiently large (e.g. $\gamma=A$ from \autoref{eqn:Diophantine_repulsion} works), such that the following holds: if for some $t',t''\in[0,T_i]$ we have $d(g_{t'}x_i,g_{t''}x_i)\leq 10r$ then $|t'-t''|\leq 10r$ (where the distance is the Euclidean one in any of the fixed rectifying coordinate charts).
If a $t''$ outside this interval existed, it would contradict the Diophantine lower bound on recurrence from \autoref{eqn:Diophantine_repulsion} for $g_{t'}x_i$.
The same property holds for $y_{i}$.

We claim that there exists $x':=g_{t'}x_i$ with $t'\in [0, T_i]$ such that the ball of radius $r$ around $x'$ (in any of the fixed coordinate charts and Euclidean metric) is disjoint from the orbit segment $[y_i,g_{T_i}y_i]$.
Assuming this first, we derive a contradiction as follows.

We pair the distributions in \autoref{lemma:Endpoint-distributions} against the bump function $\rho_{ r, x'}$:
\[
	\ip{\delta_{[x_i, g_{T_i}x_i]},\rho_{ r,x'}} \geqapprox r
	\text{ and }
	\ip{\delta_{[y_i, g_{T_i}y_i]},\rho_{ r,x'}} =0.
\]
At the same time, using the estimate in \autoref{eqn:Lipschitz-bound-xi} together with the exponential closeness of $x_i, y_i$ and $g_{T_i}x_i,g_{T_i}y_i$, we find:
\begin{align*}
	\norm{\delta_{[x_i,g_{T_i}x_i]} - \delta_{[y_i,g_{T_i}y_i]}}_{\cV^{-k_0}}
	= 
	\norm{\left(\xi_{g_{T_i}x_i} - \xi_{g_{T_i}y_i}\right)
	-
	\left(\xi_{x_i} - \xi_{y_i}\right)}_{\cV^{-k_0}}
	\leqapprox
	e^{-\ve T_i}
\end{align*}
Since $\norm{\rho_{ r,x'}}_{\cC^{k_0}}\leqapprox r^{-k_0}\asymp T_i^{k_0\gamma}$ a contradiction follows for $T_i$ large enough:
\begin{align*}
	 bT_i^{-\gamma} & = r 
	\leqapprox \ip{\delta_{[x_i, g_{T_i}x_i]},\rho_{ r,x'}}
	\\
	& \leq 
	\abs{\ip{\delta_{[y_i, g_{T_i}y_i]},\rho_{ r,x'}}}
	+
	\norm{\delta_{[x_i,g_{T_i}x_i]} - \delta_{[y_i,g_{T_i}y_i]}}_{\cV^{-k_0}}
	\cdot
	\norm{\rho_{ r,x'}}_{\cC^{k_0}}\\
	& \leqapprox
	e^{-\ve T_i}T_i^{k_0 \gamma}.
\end{align*}
It remains to show such an $x':=g_{t'}x_i$ exists.
Suppose it did not, in other words we have $B(g_{t'}x_i, r)$ intersects the orbit segment $[y_i, g_{T_i}y_i]$ for every $t'\in [0,T_i]$.
Now the existence of a reparametrization $\rho\colon [0,T_i]\to [0,T_i]$ such that
\[
	d(g_{\rho(t)} x_i, g_ty_i) \leqapprox r \asymp T_i^{-\gamma}
\]
follows from the general purpose \autoref{lem:orbit_matching}, and this contradicts \autoref{lemma:Exponential-recurrence-flows}\autoref{item:exp_rec_separation_flows}.
\hfill \qed


\begin{lemma}[Orbit matching]
	\label{lem:orbit_matching}
	Suppose $X$ is a smooth vector field with no fixed points on a compact manifold $M$, with flow $g_t$.
	Assume that for two points $x,y\in M$ and time $T>0$, their orbit segments have the following properties, for some $r>0$, $C\geq 1$:
	\begin{enumerate}
		\item If $d(g_{t'}x,g_{t''}x)<10r$ for some $t',t''\in [0,T]$, then $|t'-t''|<Cr$, and same for $y$.
		\item The orbit segment $[x,g_Tx]$ is contained in the $r$-neighborhood of $[y,g_Ty]$.
		\item We have $d(x,y)<r$ and $d(g_Tx,g_Ty)<r$.
	\end{enumerate}

	Then, provided $r$ is sufficiently small depending on $C$, there exists an orientation-preserving homeomorphism $\rho\colon [0,T]\to [0,T]$ such that
	\[
		d(g_{\rho(t)}x, g_t y)<C' r
	\]
	for some $C'>0$ depending only on $M,X,C$ but not on $r$.
\end{lemma}
\begin{proof}
	Fix a smooth hyperplane field $H(x)\subset TM(x)$ inside the tangent bundle of $M$, transverse to the vector field $X$, for example $H=X^\perp$ for the background Riemannian metric.
	Denote by $H(S,R)$ the restriction of the vectors of length strictly less than $R$ in the bundle $H$ to a subset $S\subset M$.
	Since $M$ is compact, there exist uniform constants $t_H,r_H>0$ such that the exponential map is a smooth $C^2$-bounded diffeomorphism between $H((z,g_{t_H}z),r_H)$ and an open subset of $M$, uniformly in $z\in M$, and where $(z,g_{t_H} z)$ denotes the corresponding open orbit segment.

	Fix rectifying coordinate charts as in \autoref{ssec:preliminary_constructions_in_charts}.
	There exist $K_1,K_2>0$ such that using the exponential map on $H([x,g_T x],K_1 r)$, for every $t\in [0,T]$ there exists a unique $\tilde{\rho}(t)\in [-K_2r, T+K_2r]$ and $v_t\in H(g_t x)$ such that $\exp(v_t)=g_{\tilde{\rho}(t)}y$.
	Indeed, existence follows from the containment of the orbit in the $r$-neighborhood (and assuming $r$ is sufficiently small) and uniqueness follows from the no close returns of the orbit plus the local diffeomorphism property of the exponential map.
	Furthermore, the map $\tilde{\rho}$ is a local diffeomorphism by the implicit function theorem, provided $r$ is sufficiently small.
	We thus obtain an orientation-preserving homeomorphism $\tilde{\rho}\colon [0,T] \to [a, T+b]$ where $|a|,|b|\leq K_2 r$ such that
	\[
		d(g_{t}x, g_{\tilde{\rho}(t)}y)\leq K_2 r \quad \forall t\in [0,T].
	\]
	After an affine rescaling and considering the inverse, the desired $\rho$ follows.
\end{proof}



%%					end of subsec: Proof of zero entropy
%%=============================================================================


%%%%%%%%%%%%%%%%%%%%%%%%%%%%%%%%%%%%%%%%%%%%%%%%%%%%%%%%%%%%%%%%%%%%%%%%%%%%%%%
%%%%%%%%%%%% End of Section Proofs %%%%%%%%%%%%%%%%%%%%
%%%%%%%%%%%%%%%%%%%%%%%%%%%%%%%%%%%%%%%%%%%%%%%%%%%%%%%%%%%%%%%%%%%%%%%%%%%%%%%


%%%%%%%%%%%%%%%%%%%%%%%%%%%%%%%%%%%%%%%%%%%%%%%%%%%%%%%%%%%%%%%%%%%%%%%%%%%%%%%
%%%%%%%%%%%%%%%%%%%%%%%%%%%%%%%%%%%%%%%%%%%%%%%%%%%%%%%%%%%%%%%%%%%%%%%%%%%%%%%
%						End of Main Article
%%%%%%%%%%%%%%%%%%%%%%%%%%%%%%%%%%%%%%%%%%%%%%%%%%%%%%%%%%%%%%%%%%%%%%%%%%%%%%%
%%%%%%%%%%%%%%%%%%%%%%%%%%%%%%%%%%%%%%%%%%%%%%%%%%%%%%%%%%%%%%%%%%%%%%%%%%%%%%%


%%%%%%%%%%%%%%%%%%%%%%%%%%%%%%%%%%%%%%%%%%%%%%%%%%%%%%%%%%%%%%%%%%%%%%%%%%%%%%%
%						Appendix
%%%%%%%%%%%%%%%%%%%%%%%%%%%%%%%%%%%%%%%%%%%%%%%%%%%%%%%%%%%%%%%%%%%%%%%%%%%%%%%

%\appendix
%\section{}


%%%%%%%%%%%%%%%%%%%%%%%%%%%%%%%%%%%%%%%%%%%%%%%%%%%%%%%%%%%%%%%%%%%%%%%%%%%%%%%
%						End of Appendix
%%%%%%%%%%%%%%%%%%%%%%%%%%%%%%%%%%%%%%%%%%%%%%%%%%%%%%%%%%%%%%%%%%%%%%%%%%%%%%%


\providecommand{\bysame}{\leavevmode ---\ }
\providecommand{\andchar}{\&}

\begin{thebibliography}{RHRH06}

\bibitem[BK83]{BrinKatok1983_On-local-entropy}
\textsc{Brin~M.{ \andchar\hskip 0.5em}Katok~A.} --- {``On local entropy''}. in
  {\em Geometric dynamics ({R}io de {J}aneiro, 1981)} --- vol.~1007 of {\em
  Lecture Notes in Math.}, pp.~30--38. Springer, Berlin --- 1983.
  \url{https://doi.org/10.1007/BFb0061408}.

\bibitem[CC00]{ChenChi2000_Hypoelliptic-vector-fields-and-almost-periodic-motions-on-the-torus}
\textsc{Chen~W.{ \andchar\hskip 0.5em}Chi~M.~Y.} --- {``Hypoelliptic vector
  fields and almost periodic motions on the torus {$T^n$}''}. {\em Comm.
  Partial Differential Equations} {\bfseries 25} no.~1-2, (2000) 337--354.
  \url{https://doi.org/10.1080/03605300008821509}.


\bibitem[For08]{Forni2008_On-the-Greenfield-Wallach-and-Katok-conjectures-in-dimension-three}
\textsc{Forni~G.} --- {``On the {G}reenfield-{W}allach and {K}atok conjectures
  in dimension three''}. in {\em Geometric and probabilistic structures in
  dynamics} --- vol.~469 of {\em Contemp. Math.}, pp.~197--213. Amer. Math.
  Soc., Providence, RI --- 2008. \url{https://doi.org/10.1090/conm/469/09167}.

\bibitem[GW73]{GreenfieldWallach1973_Globally-hypoelliptic-vector-fields}
\textsc{Greenfield~S.~J.{ \andchar\hskip 0.5em}Wallach~N.~R.} --- {``Globally
  hypoelliptic vector fields''}. {\em Topology} {\bfseries 12} (1973) 247--254.
  \url{https://doi.org/10.1016/0040-9383(73)90011-6}.

\bibitem[Kat01]{Katok2001_Cocycles-cohomology-and-combinatorial-constructions-in-ergodic-theory}
\textsc{Katok~A.} --- {``Cocycles, cohomology and combinatorial constructions
  in ergodic theory''}. in {\em Smooth ergodic theory and its applications
  ({S}eattle, {WA}, 1999)} --- vol.~69 of {\em Proc. Sympos. Pure Math.},
  pp.~107--173. Amer. Math. Soc., Providence, RI --- 2001.
  \url{https://doi.org/10.1090/pspum/069/1858535}.

\bibitem[RHRH06]{Rodriguez-HertzRodriguez-Hertz2006_Cohomology-free-systems-and-the-first-Betti-number}
\textsc{Rodriguez~Hertz~F.{ \andchar\hskip 0.5em}Rodriguez~Hertz~J.} ---
  {``Cohomology free systems and the first {B}etti number''}. {\em Discrete
  Contin. Dyn. Syst.} {\bfseries 15} no.~1, (2006) 193--196.
  \url{https://doi.org/10.3934/dcds.2006.15.193}.

\bibitem[Tho90]{Thomas1990_Topological-entropy-of-fixed-point-free-flows}
\textsc{Thomas~R.~F.} --- {``Topological entropy of fixed-point free flows''}.
  {\em Trans. Amer. Math. Soc.} {\bfseries 319} no.~2, (1990) 601--618.
  \url{https://doi.org/10.2307/2001256}.

\end{thebibliography}
\end{document}
